\chapter{Codearchitektur für Bedingungen}
\section*{Motivation}
Alignment-Regeln müssen mathematisch bedeutsam und rechnerisch handhabbar
dargestellt werden. AIM drückt alle Regeln als Bedingungsresiduen aus, die in
einer Optimierungsschleife einheitlich bewertet werden. Ziel ist keine Sammlung
isolierter Prüfungen, sondern eine modulare Bedingungsbibliothek.

\section{Allgemeines Prinzip}
\begin{definition}[Bedingungsresiduum]
Eine Bedingung wird als Funktion
\[
c(\mathbf x)\in\R^m
\]
dargestellt, deren Wert im Idealfall verschwindet und deren Betrag die
Verletzung misst.
\end{definition}
Dasselbe mathematische Objekt dient Machbarkeitsprüfungen, Strafgliedern und
der unmittelbaren Aufnahme in SQP-Löser.
\begin{verbatim}
function evaluateConstraint(ctx) {
  return [residual1, residual2, ...]
}
\end{verbatim}

\section{Bedingungsregister}
Eine praktische Implementierung organisiert Bedingungen in einem Register.
\begin{verbatim}
const ConstraintRegistry = {
  geometric: {}, dynamic: {}, topologic: {}, regulatory: {}
}
\end{verbatim}
Dies hält die Domänenlogik erweiterbar und vermeidet fest kodierte
monolithische Löserstrukturen.

\section{Objektmodell der Bedingung}
\begin{verbatim}
const constraint = {
  id: "min_radius",
  group: "regulatory",
  appliesTo: "edge",
  weight: 1.0,
  evaluate: (ctx) => []
}
\end{verbatim}
\texttt{id} ist der eindeutige Name, \texttt{group} die semantische Klasse,
\texttt{appliesTo} bezeichnet Kante, Knoten, Paar oder Netz,
\texttt{weight} die optionale Skalierung und \texttt{evaluate} den
Residuen-Callback.

\section{Kontextobjekt}
Bedingungen sollen nicht auf globalen Roharrays arbeiten, sondern ein
normalisiertes Kontextobjekt erhalten.
\begin{verbatim}
const ctx = {
  edge, node, pair, i, kappa, phi,
  geometry, speed, params, metadata
}
\end{verbatim}
Dies stellt eine stabile Schnittstelle zwischen Modell, Löser und Regelsystem
bereit.

\section{Geometrische Bedingungen}
\subsection{Positionsstetigkeit}
\begin{verbatim}
function continuityConstraint(p1, p2) {
  return [p1.x - p2.x, p1.y - p2.y]
}
\end{verbatim}
Dies erzwingt \(C^0\)-Stetigkeit zwischen benachbarten Elementen oder
verbundenen Alignments.

\subsection{Richtungsstetigkeit}
\begin{verbatim}
function directionConstraint(t1, t2) {
  return [t1.dx - t2.dx, t1.dy - t2.dy]
}
\end{verbatim}
Dies entspricht Tangentenstetigkeit und geometrischem \(C^1\)-Verhalten.

\subsection{Krümmungsstetigkeit}
\begin{verbatim}
function curvatureConstraint(k1, k2) {
  return [k1 - k2]
}
\end{verbatim}
Dies erzwingt stetige Krümmung über Grenzen hinweg.

\section{Dynamische Bedingungen}
\subsection{Gleichgewichtsbedingung}
\begin{verbatim}
function equilibriumConstraint(kappa, phi, v, g) {
  return [v * v * kappa - g * Math.sin(phi)]
}
\end{verbatim}
Das Residuum misst Überhöhungsfehlbetrag oder -überschuss analytisch.

\subsection{Ruckbedingung}
\begin{verbatim}
function jerkConstraint(kappa0, kappa1, phi0, phi1, ds, v, g) {
  const dk = (kappa1 - kappa0) / ds
  const dphi = (phi1 - phi0) / ds
  const j = v * v * v * dk - g * v * Math.cos(phi0) * dphi
  return [j]
}
\end{verbatim}
Sie bildet dynamische Glattheit und komfortbezogene Grenzen ab.

\section{Regulatorische Bedingungen}
\subsection{Mindestradius}
\begin{verbatim}
function minRadiusConstraint(kappa, Rmin) {
  return [Math.abs(kappa) - 1 / Rmin]
}
\end{verbatim}
Ein nichtpositiver Wert zeigt die Einhaltung der Radiusgrenze.

\subsection{Maximale Überhöhung}
\begin{verbatim}
function maxCantConstraint(phi, phiMax) {
  return [Math.abs(phi) - phiMax]
}
\end{verbatim}

\subsection{Überhöhungsgradient}
\begin{verbatim}
function cantGradientConstraint(phi0, phi1, ds, limit) {
  return [Math.abs((phi1 - phi0) / ds) - limit]
}
\end{verbatim}

\section{Topologische Bedingungen}
\subsection{Weichenknotenbedingung}
\begin{verbatim}
function switchConstraint(main, branch) {
  return [main.x - branch.x, main.y - branch.y,
          main.dx - branch.dx, main.dy - branch.dy]
}
\end{verbatim}
Dies drückt die grundlegende Kopplung an Verzweigungsknoten aus.

\subsection{Parallelgleisabstand}
\begin{verbatim}
function parallelConstraint(p1, p2, d) {
  const dx = p1.x - p2.x
  const dy = p1.y - p2.y
  return [Math.sqrt(dx * dx + dy * dy) - d]
}
\end{verbatim}

\section{Zusammenstellung der Bedingungen}
Der vollständige Residuenvektor wird aus allen aktiven Bedingungen gebildet.
\begin{verbatim}
function buildConstraintResiduals(ctx, constraints) {
  const r = []
  for (const c of constraints) {
    const ri = c.evaluate(ctx)
    const w = Math.sqrt(c.weight ?? 1)
    for (const v of ri) r.push(w * v)
  }
  return r
}
\end{verbatim}
So entsteht unabhängig von der semantischen Quelle ein löserkompatibler Vektor.

\section{Factory-Funktionen}
Bedingungen sollen durch Factories statt durch verteilte anonyme Closures
erzeugt werden.
\begin{verbatim}
function makeMinRadiusConstraint(Rmin, weight = 1) {
  return {
    id: "min_radius", group: "regulatory", appliesTo: "edge", weight,
    evaluate: ({ kappa, i }) => [Math.abs(kappa[i]) - 1 / Rmin]
  }
}
\end{verbatim}
\begin{verbatim}
function makeJerkConstraint(v, g, ds, weight = 1) {
  return {
    id: "jerk", group: "dynamic", appliesTo: "edge", weight,
    evaluate: ({ kappa, phi, i }) => {
      const dk = (kappa[i + 1] - kappa[i]) / ds
      const dphi = (phi[i + 1] - phi[i]) / ds
      return [v*v*v*dk - g*v*Math.cos(phi[i])*dphi]
    }
  }
}
\end{verbatim}

\section{Bedingungsvoreinstellungen}
Verschiedene Projekte benötigen verschiedene Konfigurationen.
\begin{verbatim}
const presetHighSpeed = [
  makeMinRadiusConstraint(4000, 10),
  makeJerkConstraint(v, g, ds, 8),
  makeMaxCantConstraint(phiMax, 6)
]
\end{verbatim}
\begin{verbatim}
const presetExistingLine = [
  makeMinRadiusConstraint(800, 3),
  makeJerkConstraint(v, g, ds, 2),
  makeParallelConstraint(trackDistance, 5)
]
\end{verbatim}

\section{Integration in den Löser}
\begin{verbatim}
function buildResidualVector(ctx, constraints, objectives) {
  return [
    ...buildConstraintResiduals(ctx, constraints),
    ...buildObjectiveResiduals(ctx, objectives)
  ]
}
\end{verbatim}
Bedingungen und Ziele teilen dieselbe residuenbasierte Schnittstelle.

\section{Diskussion}
Die Architektur bietet Modularität, Erweiterbarkeit, Kompatibilität mit
Kleinste-Quadrate-Lösern und eine transparente Abbildung von Engineering-Regeln
auf Code. Statt Prüfungen zu verteilen, konzentriert die Bibliothek die
Domänenlogik in einer mathematischen Ebene.

\section{Zusammenfassung}
\begin{summary}
Die Bedingungsbibliothek überführt Eisenbahnregeln, Topologie und dynamische
Anforderungen in eine einheitliche residuenbasierte Architektur. Sie verbindet
theoretische Modellierung, Engineering-Vorschriften und rechnerische
Optimierung und macht AIM unmittelbar in Software implementierbar.
\end{summary}
